% \usepackage{enumitem}
% \usepackage{amsmath}
% \usepackage{amsthm}
\usepackage{bbm}

% Align Block
\newcommand\numberthis{\addtocounter{equation}{1}\tag{\theequation}}
\allowdisplaybreaks

\usepackage{mathtools}

% Math delimiters
\DeclarePairedDelimiter{\abs}{\lvert}{\rvert}
\DeclarePairedDelimiter{\brk}{[}{]}
\DeclarePairedDelimiter{\crl}{\{}{\}}
\DeclarePairedDelimiter{\prn}{(}{)}
\DeclarePairedDelimiter{\nrm}{\|}{\|}
\DeclarePairedDelimiter{\tri}{\langle}{\rangle}
\DeclarePairedDelimiter{\dtri}{\llangle}{\rrangle}
\DeclarePairedDelimiter{\ceil}{\lceil}{\rceil}
\DeclarePairedDelimiter{\floor}{\lfloor}{\rfloor}
\newcommand{\midsem}{\,;}
\newcommand{\trc}[1]{\text{tr} \prn*{#1}}

\let\Pr\undefined
\let\P\undefined
\DeclareMathOperator{\En}{\mathbb{E}}
\newcommand{\Ep}{\wt{\bbE}}
%\DeclareMathOperator*{\Eh}{\widehat{\mathbb{E}}}
\DeclareMathOperator{\P}{P}
\DeclareMathOperator{\Pr}{\bbP}

% Arg<x>
\DeclareMathOperator*{\argmin}{argmin} % * Places subscript directly under operator
\DeclareMathOperator*{\argmax}{argmax}
\DeclareMathOperator*{\arginf}{arginf}
\DeclareMathOperator*{\argsup}{argsup}

% one-off macros
\newcommand{\ls}{\ell}
\newcommand{\bls}{\mb{\ls}}
\newcommand{\ind}[1]{\mathbbm{1}\crl*{#1}}    %Indicator
\newcommand{\indd}[1]{\mathbbm{1}\crl{#1}}    %Indicator
\newcommand{\eps}{\varepsilon}
\newcommand{\veps}{\varepsilon}

\newcommand{\ldef}{\vcentcolon=}
\newcommand{\rdef}{=\vcentcolon}

% styles
\newcommand{\mc}[1]{\mathcal{#1}}
\newcommand{\wt}[1]{\widetilde{#1}}
\newcommand{\wh}[1]{\widehat{#1}}
%\newcommand{\mb}[1]{\boldsymbol{#1}}


% Special letters: blackboard, mathcal, widehat
% djhsu magic
\def\ddefloop#1{\ifx\ddefloop#1\else\ddef{#1}\expandafter\ddefloop\fi}
\def\ddef#1{\expandafter\def\csname bb#1\endcsname{\ensuremath{\mathbb{#1}}}}
\ddefloop ABCDEFGHIJKLMNOPQRSTUVWXYZ\ddefloop
\def\ddefloop#1{\ifx\ddefloop#1\else\ddef{#1}\expandafter\ddefloop\fi}
\def\ddef#1{\expandafter\def\csname b#1\endcsname{\ensuremath{\mathbf{#1}}}}
\ddefloop ABCDEFGHIJKLMNOPQRSTUVWXYZ\ddefloop
\def\ddef#1{\expandafter\def\csname c#1\endcsname{\ensuremath{\mathcal{#1}}}}
\ddefloop ABCDEFGHIJKLMNOPQRSTUVWXYZ\ddefloop
\def\ddef#1{\expandafter\def\csname h#1\endcsname{\ensuremath{\widehat{#1}}}}
\ddefloop ABCDEFGHIJKLMNOPQRSTUVWXYZabcdefghijklmnopqrsuvwxyz\ddefloop    % Not defined for t.
\def\ddef#1{\expandafter\def\csname hc#1\endcsname{\ensuremath{\widehat{\mathcal{#1}}}}}
\ddefloop ABCDEFGHIJKLMNOPQRSTUVWXYZ\ddefloop
\def\ddef#1{\expandafter\def\csname t#1\endcsname{\ensuremath{\widetilde{#1}}}}
\ddefloop ABCDEFGHIJKLMNOPQRSTUVWXYZ\ddefloop
\def\ddef#1{\expandafter\def\csname tc#1\endcsname{\ensuremath{\widetilde{\mathcal{#1}}}}}
\ddefloop ABCDEFGHIJKLMNOPQRSTUVWXYZ\ddefloop

% Names
\newcommand{\Holder}{H{\"o}lder}

% COMMONLY USED MACROS
%Losses
\newcommand{\KL}[2]{\mathrm{KL}{\prn*{#1 \| #2}}}
\newcommand{\kl}[2]{\mathrm{kl}{\prn*{#1 \| #2}}}
\newcommand{\ball}{\mathbb{B}}
\newcommand{\poprisk}{L_{\cD}}
\newcommand{\emprisk}{\wh{L}_n}
\newcommand{\sign}{\mathrm{sign}}

% Matrices
%\newcommand{\diag}{\textrm{diag}}
%\newcommand{\rank}{\mathrm{rank}}
\newcommand{\trn}{\dagger}
\newcommand{\Tr}{\mathbf{Tr}}

% Special Symbols
\newcommand{\xr}[1][n]{x_{1:#1}}
\newcommand{\br}[1][n]{b_{1:#1}}
\newcommand{\er}[1][n]{\eps_{1:#1}}
\newcommand{\sr}[1][n]{\sigma_{1:#1}}
\newcommand{\yr}[1][n]{y_{1:#1}}
\newcommand{\zr}[1][n]{z_{1:#1}}
\renewcommand{\th}{{\text{th}}}
%Circled Numbers
\usepackage{tikz}
\newcommand*\circled[1]{\tikz[baseline=(char.base)]{
            \node[shape=circle,draw,inner sep=2pt] (char) {#1};}}

% Calculus
\newcommand{\grad}{\nabla}

\newcommand{\proman}[1]{\prn*{\romannumeral #1}}
\newcommand{\overles}[1]{\overset{ #1}{\lesssim{}}}
\newcommand{\overleq}[1]{\overset{ #1}{\leq{}}}
\newcommand{\overgeq}[1]{\overset{#1}{\geq{}}}
\newcommand{\overeq}[1]{\overset{#1}{=}}
\newcommand{\overapx}[1]{\overset{#1}{\approx}}

%%% theorems

%\theoremstyle{definition}  %Sets style of subsequent newtheorems to 'definition'
%\newtheorem{assumption}{Assumption}
\newtheorem{condition}{Condition}
\newtheorem{claim}{Claim}
%\newtheorem{exercise}{Exercise}
\newtheorem{assumption}{Assumption}
\newtheorem{conjecture}{Conjecture}
\newtheorem{corollary}{Corollary}
\newtheorem{proposition}{Proposition}
%\newtheorem{fact}{Fact}
%\newtheorem{goal}{Goal}
%%\newtheorem{model}{Model}
\newtheorem{problem}{Problem}
%\newtheorem{model}{Model}
%\theoremstyle{plain}
{\newtheorem{lemma}{Lemma}}
\newtheorem{remark}{Remark}
%\newtheorem{example}{Example}
{\newtheorem{theorem}{Theorem}}
\newtheorem{theorem*}{Theorem}
\newtheorem{definition}{Definition}
%\newtheorem{question}{Question}
% \newtheorem{proof}{Proof}



% \newtheorem*{assumption*}{\assumptionnumber}
% \providecommand{\assumptionnumber}{}
% \makeatletter
% \newenvironment{spassumption}[1]
%  {%
%   \renewcommand{\assumptionnumber}{Assumption #1}%
%   \begin{assumption*}%
%   \protected@edef\@currentlabel{#1}%
%  }
%  {%
%   \end{assumption*}
%  }
% \makeatother

% \usepackage{xpatch}
% \xpatchcmd{\proof}{\itshape}{\normalfont\proofnameformat}{}{}
% \newcommand{\proofnameformat}{\bfseries}

% The template's prettyref helpers are unused by this method note.
% They are omitted from this copy because prettyref.sty is not available
% in the installed TinyTeX distribution.

\newcommand{\algcomment}[1]{\textcolor{blue!70!black}{\footnotesize{\texttt{\textbf{//
          #1}}}}}
\newcommand{\algcommentbig}[1]{\textcolor{blue!70!black}{\footnotesize{\texttt{\textbf{/*
          #1~*/}}}}}
